\begin{table}[htbp]
\centering\small
\caption{Matched directed tolerance on a new finite query distribution}\label{tab:r19warm}
\begin{tabular}{rrrrrr}\toprule
$d$ & Target & Neural GD & Zero GD & Accelerated & Wins / 96\\
\midrule
8 & $10^{-2}$ & 1.05 & 3.44 & 3.84 & 96\\
8 & $10^{-3}$ & 2.35 & 5.19 & 5.12 & 96\\
8 & $10^{-4}$ & 4.27 & 7.16 & 6.97 & 96\\
8 & $10^{-6}$ & 8.56 & 11.56 & 10.84 & 96\\
32 & $10^{-2}$ & 1.86 & 4.72 & 4.94 & 96\\
32 & $10^{-3}$ & 3.31 & 6.47 & 6.06 & 96\\
32 & $10^{-4}$ & 5.36 & 8.66 & 8.00 & 96\\
32 & $10^{-6}$ & 9.84 & 13.00 & 12.00 & 96\\
128 & $10^{-2}$ & 5.00 & 6.00 & 5.97 & 93\\
128 & $10^{-3}$ & 7.00 & 8.00 & 7.00 & 0\\
128 & $10^{-4}$ & 9.06 & 10.00 & 9.00 & 0\\
128 & $10^{-6}$ & 14.00 & 14.00 & 13.00 & 0\\
\bottomrule
\end{tabular}
\par\smallskip\begin{minipage}{.98\textwidth}\footnotesize Mean correction counts; each method uses one gradient evaluation per correction and one final gradient evaluation. Neural means pool three frozen initializers and 32 predeclared queries. Classical means use the same 32 queries. A win means strictly fewer gradient evaluations than acceleration for the paired query; it is not a wall-time claim. Acceleration returns its independently checked search point. The set of 32 queries is not a cover of the state cube.\end{minipage}
\end{table}
